% =====================================================================
% Supervisor meeting, 20 August 2026
% "The Iberian setting: what the data holds, and what killed each design"
%
% Build:  bash slides/build-tex.sh
%   or:   cd slides && xelatex supervisor-2026-08-19-data-and-designs.tex
%
% Compile from INSIDE slides/. Running xelatex from the repo root with
% -output-directory=slides fails to load the figures: kpathsea resolves
% "figures/" through the output directory and reports a hit, but the graphics
% loader then reads the path relative to the working directory and finds
% nothing. It exits 0 with a 2-page PDF, so the failure is easy to miss.
%
% Figures come from scripts/analysis/slides_data_figures.py and
% scripts/analysis/slides_figures.py -- regenerate with:
%   PYTHONPATH=scripts/analysis .venv/bin/python scripts/analysis/slides_data_figures.py
% =====================================================================
\documentclass[9pt,aspectratio=169]{beamer}

% Metropolis, tuned light: the theme ships a dark bar behind the frame title;
% this deck uses a plain title with a thin accent rule under it instead, so the
% figures (which are drawn on white) sit on the same background as the slide.
\usetheme[titleformat frame=smallcaps,
          progressbar=frametitle,
          numbering=fraction,
          block=fill]{metropolis}

% ---------- fonts (xelatex) ------------------------------------------
\usepackage{fontspec}
\setmainfont{Helvetica Neue}
\setsansfont{Helvetica Neue}
\setmonofont{Menlo}[Scale=0.9]
\usepackage{amsmath}          % load before unicode-math
\usepackage{unicode-math}
\setmathfont{STIX Two Math}

% ---------- packages --------------------------------------------------
\usepackage{booktabs}
\usepackage{tabularx}
\usepackage{graphicx}
\usepackage{etoolbox}
\graphicspath{{figures/}{slides/figures/}}

% Unicode the sans font lacks, mapped to math equivalents, so you can type
% readable characters straight into the source instead of LaTeX escapes.
\usepackage{newunicodechar}
\newunicodechar{→}{\ensuremath{\rightarrow}}
\newunicodechar{×}{\ensuremath{\times}}
\newunicodechar{≈}{\ensuremath{\approx}}
\newunicodechar{−}{\ensuremath{-}}
\newunicodechar{β}{\ensuremath{\beta}}
\newunicodechar{η}{\ensuremath{\eta}}
\newunicodechar{Δ}{\ensuremath{\Delta}}
% Helvetica Neue carries the euro glyph, so \euro needs no extra package.
\providecommand{\euro}{€}

% ---------- look ------------------------------------------------------
% Accent is the blue the figures already use (dataviz palette), so the rule
% under each title and the chart lines are the same colour.
\definecolor{accent}{HTML}{2A78D6}
\definecolor{ink}{HTML}{0B0B0B}
\definecolor{ink2}{HTML}{52514E}

\setbeamercolor{normal text}{fg=ink,bg=white}
\setbeamercolor{frametitle}{fg=ink,bg=white}      % kill Metropolis' dark bar
\setbeamercolor{progress bar}{fg=accent,bg=accent!25}
\setbeamercolor{alerted text}{fg=accent}
\setbeamerfont{frametitle}{size=\large,series=\bfseries}
\setbeamerfont{page number in head/foot}{size=\scriptsize}
\setbeamercolor{page number in head/foot}{fg=ink2}

% Do NOT override \setbeamertemplate{frametitle} here: Metropolis draws the
% smallcaps title AND the accent progress bar from that template, so replacing
% it silently removes both. Adjust padding with the theme's own length instead.
\makeatletter
\setlength{\metropolis@frametitle@padding}{1.4ex}
\makeatother

\setlength{\aboverulesep}{0pt}
\setlength{\belowrulesep}{0pt}
\renewcommand{\arraystretch}{1.25}

% Every slide's table is set at \small and full text width. Two shapes are
% used throughout, so changing these two macros restyles the whole deck.
%
%   \designtable{<left column width>}{ rows... }   label + free text
%   \figframe{<title>}{<file>}{ caption... }       one figure, one caption
%
\newcommand{\designtable}[2]{%
  \small
  \begin{tabularx}{\textwidth}{@{}>{\raggedright\arraybackslash}p{#1}X@{}}
  \toprule
  #2
  \bottomrule
  \end{tabularx}%
}
\newcommand{\figframe}[3]{%
  \begin{frame}{#1}
    \centering\includegraphics[width=\textwidth]{#2}\par
    \vspace{0.6em}
    \raggedright\small #3
  \end{frame}%
}

\title{The Iberian setting: what the data holds, and what killed each design}
\author{Theodor Mello}
\date{20 August 2026}

\begin{document}

\begin{frame}[plain]
  \titlepage
\end{frame}

% =====================================================================
% 1. The decision
% =====================================================================
\begin{frame}{The decision on the table}
  \textbf{Keep the Iberian setting, or kill it.}

  \vspace{0.8em}
  Three things to put to you:

  \vspace{0.4em}
  \begin{enumerate}
    \item \textbf{The data} is good --- and it is built, audited and
          documented.
    \item \textbf{One result} already holds up: market power roughly doubled
          while measured concentration \emph{fell}.
    \item \textbf{Twelve designs have died.} Not twelve accidents --- three
          recurring causes, and they are properties of the setting.
  \end{enumerate}

\end{frame}

% =====================================================================
% 2. What we hold
% =====================================================================
\figframe{What we hold: nine years, fully re-fetchable}
         {data_coverage.png}
         {Both Spanish interventions sit \textbf{inside} the window:
          \textbf{3.7 years before the levy}, and \textbf{2.4 years after the
          cap ended}.}

% =====================================================================
% 3. Why the data is unusual
% =====================================================================
\begin{frame}{Why this data is unusual}
  \designtable{0.24\textwidth}{
    \textbf{Full bid curves} &
    354\,M unit-level day-ahead offers --- \textbf{accepted \emph{and}
    rejected}. We observe firms' supply \emph{functions}, not just realised
    supply. \\
    \midrule
    \textbf{Declared start-up costs} &
    Every plant states, each day, the minimum revenue it needs to run --- a
    fixed \euro/day and a variable \euro/MWh term. \\
    \midrule
    \textbf{Observable contract position} &
    The volume sold forward, which blunts the incentive to withhold. Bilateral contracts run
    \textbf{39--49\%} of scheduled sell volume. \\
    \midrule
    \textbf{Two countries, one auction} &
    Spain and Portugal clear jointly against one merged bid stack --- same
    price, same forward curve, different regulator. \\
  }
\end{frame}

% =====================================================================
% 4. The window
% =====================================================================
\figframe{The window: a price shock, and a VRE share that doubled}
         {price_vre.png}
         {Day-ahead price 34 \euro/MWh (2020) → 168 (2022) → 46 (2026).
          Wind and solar 24\% → 42\% of Spanish generation.}

% =====================================================================
% 5. Who sets the price
% =====================================================================
\figframe{Who sets the price}
         {marginal_tech.png}
         {\textbf{Hydro is the modal marginal technology in every year} ---
          30--52\% of periods. Gas is second and collapses after the cap. This
          is why every design here ended up about hydro.}

% =====================================================================
% 6. Ownership mapping
% =====================================================================
\begin{frame}{Who owns what: the mapping, and its ceiling}
  \small
  OMIE publishes no owner. Built in four layers from ESIOS, \texttt{cab} plant
  names and a hand-verified prefix dictionary.

  \vspace{0.8em}
  \centering
  \begin{tabular}{@{}lr@{}}
    \toprule
    what & coverage \\
    \midrule
    Units in the registry (ES 4,960 / PT 150)            & 5,137 \\
    \textbf{Near-marginal offered volume → market agent} & \textbf{89.1\%} \\
    Generation volume → parent firm (68 firms)           & 82.4\% \\
    \textbf{Dispatchable capacity → parent firm}         & \textbf{96.0\%} \\
    Generation volume bid by a \emph{named plant} (technology readable) & 47.7\% \\
    \bottomrule
  \end{tabular}

  \vspace{1.0em}
  \raggedright
  \textbf{Why only 48\%?} The rest is \textbf{not unparsed --- it is bid by
  portfolios, not plants}. \texttt{EDPSVD1} ``VENTA RE'' is 6,500 MW of mixed
  renewables in one unit: there is no single technology to name.
\end{frame}

% =====================================================================
% 7. The result
% =====================================================================
\figframe{The result that holds: market power doubled as concentration fell}
         {concentration_markup.png}
         {$L_i = s_i/\eta$. Shares barely moved; \textbf{$\eta$ halved} ---
          1.83 (2018) → 0.85 (2023). For Iberdrola/Endesa the implied Lerner
          floor goes \textbf{11\% → 21\%}, and to \textbf{34\%} in the
          highest-VRE third of 2023 hours.}

% =====================================================================
% 8. The mechanism
% =====================================================================
\figframe{The mechanism: renewables thin the market \emph{at the margin}}
         {eta_vre.png}
         {Wind and solar bid at zero and sit far inframarginal --- they do not
          flatten the curve where the price is set, they push the crossing
          point into a region with fewer flexible units.}

% =====================================================================
% 11. The diagnosis
% =====================================================================
\begin{frame}{The diagnosis: three causes, not twelve accidents}
  \designtable{0.04\textwidth}{
    \textbf{1} &
    \textbf{The treatment has no cross-section.} Levy, gas cap and renewable
    share hit every unit at once. Date fixed effects and the design are
    mutually exclusive. \\
    \textbf{2} &
    \textbf{Any summary of ``how a plant bid'' changes sign depending on how
    it is built.} The median sell offer is \euro 0 and the 99th percentile
    \euro 1,000, so a few extreme offers decide the answer --- and none of them
    can be read as withholding rather than a genuinely high cost. \\
    \textbf{3} &
    \textbf{A hard cluster ceiling.} Only 32 parent firms own dispatchable
    plant. Any treatment varying at firm level has a few dozen clusters. \\
    \textbf{+3} &
    \textbf{No usable comparison, or too few observations.} The single-date
    event study, the run-of-river control group and the 10 MW RDD each failed
    before estimation. \\
  }

  \vspace{0.9em}
  \small
  \textbf{Everything that is a \emph{direct computation} held}: the elasticity
  decline, firm and technology coverage, contract shares, firm-specific
  residual demand, the implied-marginal-cost ordering by technology.

  \vspace{0.4em}
  Everything built as a \emph{constructed behavioural measure} collapsed.
\end{frame}


% =====================================================================
% APPENDIX -- derivations. \appendix keeps these out of the slide count.
% =====================================================================
\appendix

\begin{frame}[plain]
  \centering\vspace{2em}
  {\Large\bfseries Appendix: how the three numbers are built}
\end{frame}

% ---------------------------------------------------------------------
% A1. Residual demand elasticity
% ---------------------------------------------------------------------
\begin{frame}{A1. Residual demand, and its elasticity $\eta$}
  \small
  Firm $i$ does not face market demand, it faces what is left after every
  rival has supplied. For period $t$:
  \[
    RD_{it}(p) \;=\; D_t(p) \;-\; \sum_{j \neq i} S_{jt}(p)
  \]
  We build the \emph{market-wide} version, $RD_t(p) = D_t(p) - S_t(p)$, which
  is the curve facing an infinitesimally small participant. Its slope at the
  clearing price $p^*_t$ is
  \[
    \bigl| RD'_t(p^*_t) \bigr| \;=\; \bigl| D'_t(p^*_t) \bigr| \;+\; S'_t(p^*_t)
  \]
  Submitted curves are \textbf{step functions and have no derivative}, so the
  slope is estimated as the density of offered volume in a price band of width
  $2h$ around the clearing price:
  \[
    \bigl| RD'_t(p^*_t) \bigr| \;\approx\;
      \frac{\text{sell MW} + \text{buy MW within } p^*_t \pm h}{2h}
    \qquad\Longrightarrow\qquad
    \eta_t \;=\; \frac{\bigl| RD'_t \bigr| \, p^*_t}{Q_t}
  \]
  $\eta_t$ is scale-free, so it is comparable across years in which the price
  level moved by a factor of five.
\end{frame}

\begin{frame}{A1. How $\eta$ is applied to the data}
  \designtable{0.30\textwidth}{
    Unit of observation & One market period (one settlement interval for the
    whole Iberian market) \\
    Frequency & Hourly to 30 September 2025; \textbf{quarter-hourly} from
    1 October 2025, when the market changed its settlement interval \\
    Sample & \textbf{79,036 periods}, 2018--2026 --- 62,721 hourly and
    16,315 quarter-hourly \\
    Input & Every \emph{submitted} bid and offer, accepted \textbf{and}
    rejected (354 M rows). Cleared-only data cannot give a slope \\
    Bandwidth $h$ & \textbf{5 \euro/MWh} in the reported results; 1, 2 and 10
    also computed, rank correlation 0.92--0.94 \\
    Units & $|RD'|$ in MW per \euro/MWh; $p^*$ in \euro/MWh; $Q$ in MWh;
    $\eta$ dimensionless \\
    Periods dropped & Market splitting (Spain and Portugal price apart);
    whole-day plant dropouts; empty price bands; $p^* \leq 1$ \euro/MWh,
    where $\eta$ is not defined \\
  }
\end{frame}

% ---------------------------------------------------------------------
% A2. The markup
% ---------------------------------------------------------------------
\begin{frame}{A2. From $\eta$ to the markup}
  \small
  Firm $i$ picks its output taking the residual demand curve as given:
  \[
    \max_{q_i} \; p(q_i)\, q_i - C_i(q_i)
    \qquad\Longrightarrow\qquad
    p + q_i \frac{\partial p}{\partial q_i} = C_i'(q_i)
  \]
  and since $\partial p / \partial q_i = 1 / RD_i'(p)$, the markup over
  marginal cost is
  \[
    p - C_i'
    \;=\; \frac{q_i}{\bigl| RD_i' \bigr|}
    \qquad\Longrightarrow\qquad
    L_i \;\equiv\; \frac{p - C_i'}{p}
    \;=\; \frac{q_i}{\bigl| RD_i' \bigr| \, p}
    \;=\; \frac{s_i}{\eta}
  \]
  with $s_i = q_i / Q$. Volume already sold forward at a fixed price earns
  nothing extra when $p$ rises, so with a contract position $\theta_i$ the
  markup becomes $\bigl(q_i - \theta_i\bigr) / |RD_i'|$.

  \vspace{0.5em}
  \textbf{Why this is a lower bound.} A firm's own supply is excluded from its
  own residual demand, so $|RD_i'| \leq |RD'|$: every real firm faces a curve
  \emph{at least} as steep as the market-wide one. Hence
  $L_i \geq s_i / \eta$ --- the least markup consistent with profit
  maximisation, not an estimate of the markup taken.
\end{frame}

\begin{frame}{A2. How the markup is applied to the data}
  \designtable{0.30\textwidth}{
    $\eta$ & Median across all periods within the year, from A1 \\
    $s_i$ & Annual generation share by business group, published by the
    competition regulator, covering Spain and Portugal together, 2018--2023 ---
    so the yearly figures pair an annual share with a per-period elasticity \\
    Also computed & Share per firm \emph{per period} from cleared volume,
    1,101,980 firm-period observations, which avoids treating a share as
    constant across a year (within-year variation is large) \\
    Units & $L_i$ is a share of the price, dimensionless; the markup in
    \euro/MWh is $L_i \times p^*$ \\
    Reported & For the average of the two largest firms: $L$ rises
    \textbf{11\% (2018) → 21\% (2023)}, and to \textbf{34\%} in the
    highest-renewable third of 2023 hours \\
    Two biases, opposite & \textbf{Understated} on the elasticity side (the
    market-wide curve is flatter than any firm's own). \textbf{Overstated} on
    the share side: the published share counts volume sold bilaterally, which
    is not exposed to the day-ahead price \\
  }
\end{frame}

% ---------------------------------------------------------------------
% A3. Concentration
% ---------------------------------------------------------------------
\begin{frame}{A3. Concentration, and why it moves the other way}
  \small
  The standard concentration measure is the sum of squared shares:
  \[
    H \;=\; 10{,}000 \times \sum_i s_i^{\,2}
  \]
  \vspace{-0.6em}
  \designtable{0.30\textwidth}{
    Source & Published by the competition regulator, annual, Spain and
    Portugal together, 2008--2023 \\
    Computed on & The \emph{operating schedule} --- day-ahead cleared volume
    \textbf{plus} bilateral contracts and self-scheduling. On January 2024
    that is \textbf{1.73$\times$} day-ahead cleared volume \\
    Units & Index on 0--10,000; 1,000 is roughly ten equal firms \\
  }

  \vspace{0.8em}
  \textbf{The two measures can move in opposite directions because they share
  only one of their two ingredients.} $H$ depends on shares alone. The Lerner
  index $L_i = s_i / \eta$ depends on shares \emph{and} on the slope of
  residual demand. Between 2018 and 2023 shares barely moved --- $H$ fell
  1,314 → 1,007 --- while $\eta$ fell 1.83 → 0.85. The same share converts
  into twice the markup.
\end{frame}

\end{document}
